\section{Sáu lần chẩn đoán, hai lần đề nghị}
\label{sec:ch18-chan-doan-va-de-nghi}

\index{khoảng trống nghiên cứu!chẩn đoán và đề nghị là hai câu hỏi|(}%
Cái hình mà mục trước dựng ra có ba phần, và phần thứ ba là khoảng cách giữa điều
các bài chẩn đoán với điều các bài đề nghị. Mục này đo khoảng cách ấy, và cách đo
là hỏi lại cả chín bài một câu mà không chương nào trong bảy chương Phần~IV đã
hỏi.

\index{khoảng trống nghiên cứu!hai câu hỏi khác nhau}%
Hai câu hỏi ấy phải tách cho rõ, vì chúng cho hai con số khác nhau và cả hai đều
đúng. Câu thứ nhất là câu mà cái sổ được lập để trả lời: bài này có \emph{nói} ở
đâu đó rằng thiếu một dụng cụ đo hay không. Câu ấy nhận câu trả lời từ bất cứ chỗ
nào trong bài, một dòng kết quả, một con số đếm được, một câu giữa phần thân. Câu
thứ hai hẹp hơn: phần hạn chế hoặc phần công việc tương lai \emph{của chính bài
ấy} có đòi cái dụng cụ ấy hay không. Chỉ phần công việc tương lai trả lời được
câu thứ hai, vì đó là chỗ một bài nói ra việc nó cho là đáng làm tiếp.

Câu thứ nhất cho sáu trên chín, và mục trước đã đếm. Câu thứ hai cho hai trên
chín. Bảng~\ref{tab:ch18-so-chan-doan} đặt hai câu trả lời cạnh nhau cho từng
bài, đánh số bài theo vị trí trong thang đọc: bài 21 là bài neo mà
chương~\ref{ch:chi-so-khong-do-do-trung-thuc} đọc, các bài 22 tới 26 lần lượt
thuộc các chương~\ref{ch:lime-nao-dang-tin}, \ref{ch:tro-choi-attribution},
\ref{ch:con-nguoi-vang-mat} (hai bài) và \ref{ch:gioi-han-ly-thuyet}, còn ba bài
27, 28 và 29 đều thuộc chương~\ref{ch:giai-thich-duoi-tan-cong}.

\begin{table}[htbp]
  \centering
  \footnotesize
  \caption{Chín bài trong sổ của Phần~IV, đánh số theo vị trí trong thang đọc mà
  phụ lục~\ref{app:danh-muc} ghi cùng chương đọc từng bài. Cột giữa là câu hỏi mà
  cái sổ được lập để trả lời; cột phải là câu hỏi mà chỉ phần công việc tương lai
  trả lời được. Khoảng cách giữa hai cột là điều mục này đọc.}
  \label{tab:ch18-so-chan-doan}
  \begin{tabular}{@{}lp{4cm}p{6cm}@{}}
    \toprule
    Bài & Có nêu chỗ thiếu không & Công việc tương lai đòi gì \\
    \midrule
    21 & Có, bằng phép đo & Chỉ số đáng tin và hiệu quả hơn; chỉ nêu loại \\
    22 & Có, không có chuẩn dùng & Khung đánh giá riêng cho LIME \\
    23 & Có, bước chưa làm & \textbf{Đường cong deletion và insertion trên từng
    cặp token} \\
    24 & Có, bằng phép đếm & Tra cứu rộng hơn; tiêu chí chấm khác \\
    25 & Có, chạy một lần & Chỉ số đặt nền trên cảm nhận của người \\
    26 & Không & Một đại lượng tính được; một thí nghiệm \\
    27 & Có, về chính nó & \tn{Bộ phát hiện ngoài phân phối}{out-of-distribution
    detector}; thêm bộ dữ liệu \\
    28 & Không & Mô hình khái quát hóa tốt hơn \\
    29 & Không & Khám phá nhân quả tốt hơn \\
    \bottomrule
  \end{tabular}
\end{table}

\index{khoảng trống nghiên cứu!một câu duy nhất gọi tên dụng cụ}%
Đọc cột phải từ trên xuống thì thấy một câu duy nhất gọi tên một dụng cụ. Nó nằm
trong phần hạn chế của bài mà chương~\ref{ch:tro-choi-attribution} đọc, và
mục~\ref{sec:ch11-dieu-chua-giai-quyet} đã dẫn nó: một phép so nghiêm túc giữa
các phương pháp sẽ cần tới phép đo độ trung thực, và dụng cụ mà các tác giả sẽ
dùng là các đường cong deletion và insertion trên từng cặp
token~\autocite{p23metagame}. Tám câu còn lại không gọi tên dụng cụ nào. Bài 21
đòi các chỉ số độ trung thực đáng tin và hiệu quả hơn,%
\footnote{Nguyên văn: \enquote{more reliable and efficient faithfulness
metrics}.}
tức đòi một loại chứ không một thiết kế, và phần hạn chế riêng của nó thì bàn
chuyện khác hẳn, là phạm vi mà cách gán nhãn của nó với
tới~\autocite{p21metrics}.

\index{khoảng trống nghiên cứu!phần lớn đòi thứ mình vốn đã dựng}%
Bảy dòng còn lại của cột phải phần lớn có chung một hình dạng, và hình dạng ấy là
kết quả của mục này: bài đòi một bản tốt hơn của đúng thứ nó vốn đã dựng. Bài về
biến thể LIME đòi một khung đánh giá cho LIME. Bài dùng khám phá nhân quả đòi
thuật toán khám phá nhân quả tốt hơn. Bài đếm số nghiên cứu có người đòi một cuộc
đếm rộng hơn. Bài chứng minh cái trần đòi một đại lượng tính được thay cho đại
lượng không tính được mà cái trần của nó dựng trên. Ngay cả bài tự nhận chưa đo
xem công cụ mình có giải thích được gì cũng không đòi một dụng cụ đo: nó đòi chạy
trên nhiều bộ dữ liệu hơn~\autocite{p27shlime}, mà chạy trên nhiều bộ dữ liệu hơn
thì vẫn phải chấm bằng một cái gì, và bài không nói bằng cái gì.

Hai dòng không khớp hình dạng ấy, và chúng đáng nêu ra chứ không đáng bỏ qua. Bài
đo tỉ lệ phân loại sai đòi những mô hình khái quát hóa tốt hơn, tức đòi sửa thứ
bị tấn công chứ không sửa công cụ tấn công mà nó dựng. Và bài chạy phép đối chiếu
với người đòi một loại chỉ số mới, đặt nền trên cảm nhận của người và được kiểm
định theo cảm nhận ấy, tức đòi một dụng cụ mà nó chưa dựng. Dòng thứ hai này là
dòng gần nhất với việc đòi một dụng cụ đo trong cả cột, và nó vẫn không phải đòi một dụng cụ đo độ trung
thực: tiêu chí nó đề nghị là người đọc, tức dòng \tn{tính hợp lý}{plausibility}
của bảng~\ref{tab:ch01-bon-quan-he} chứ không phải dòng độ trung thực.

Chỗ ấy đáng dừng lại, vì nó là cả khoảng trống thu vào một câu. Một bài nhận ra
rằng nó chưa kiểm xem lời giải thích của mình có đúng không, rồi đề nghị kiểm
bằng cách làm nhiều hơn cùng một việc, là một bài chưa nhận ra rằng phép kiểm ấy
cần một dụng cụ mà chưa ai có.

\index{khoảng trống nghiên cứu!chẩn đoán và đề nghị là hai câu hỏi|)}%
Vậy khoảng cách giữa sáu và hai là khoảng cách giữa chẩn đoán và đơn thuốc. Lĩnh
vực nói ra chỗ thiếu sáu lần, từ sáu hướng, trong sáu bài của ba nhánh khác nhau.
Nó đề nghị làm gì về chỗ thiếu ấy hai lần, và một trong hai lần chỉ nêu tên loại.
Cái sổ đo được vế thứ nhất vì nó được lập để đo vế ấy; vế thứ hai phải hỏi riêng,
và mục sau hỏi nó của ba bài đứng ngoài danh mục, ba bài mà nếu ai đã dựng dụng
cụ ấy thì phải là họ.
